% SEGMENT OLP-0377-S01
% Part: lambda-calculus
% Chapter: lambda-definability
% Section: truth-values

\documentclass[../../../include/open-logic-section]{subfiles}

\begin{document}

\olfileid{lam}{ldf}{tvr}
\olsection{மெய்மதிப்புகளும் தொடர்புகளும்}

தூய லாம்டா கலனத்தில் மெய்மதிப்புகளைப் பின்வருமாறு குறியாக்கலாம்:
\begin{align*}
  \fn{true} & \ident \lambd[x][\lambd[y][x]]\\
  \fn{false} & \ident \lambd[x][\lambd[y][y]]
\end{align*}

மெய்மதிப்புகள் \emph{தேர்விகளாக}, அதாவது இரண்டு உள்ளீடுகளை ஏற்று அவற்றில்
ஒன்றைத் திருப்பித்தரும் சார்புகளாக, பிரதிநிதித்துவப்படுத்தப்படுகின்றன.
மெய்மதிப்பு $\fn{true}$ தனது முதல் உள்ளீட்டைத் தேர்ந்தெடுக்கிறது;
$\fn{false}$ தனது இரண்டாம் உள்ளீட்டைத் தேர்ந்தெடுக்கிறது. எடுத்துக்காட்டாக,
$\fn{true}\, M N$ எப்போதும் $M$ ஆகக் குறைகிறது; அதேபோல்
$\fn{false}\, M N$ எப்போதும்~$N$ ஆகக் குறைகிறது.

\begin{defn}
% SOURCE CORRECTION TA-LAM-037: the displayed arguments and quantified relation use k places, so the ambient power is Nat^k.
$R \subseteq \Nat^k$ என்ற ஒரு தொடர்புக்கு, பின்வருமாறு அமையும் சொல்~$R$
ஒன்று இருந்தால் அத்தொடர்பை !!{lambda definable} என்று அழைக்கிறோம்:
\begin{align*}
  R\, \num{n_1} \dots \num{n_k} & \bred \fn{true}
  \intertext{$R(n_1, \dots, n_k)$ அமையும் போதெல்லாம்; மேலும்}
  R\, \num{n_1} \dots \num{n_k} & \bred \fn{false}
\end{align*}
மற்ற எல்லா நிலைகளிலும்.
\end{defn}

% SOURCE CORRECTION TA-LAM-038: IsZero holds exactly of 0; remove the frozen duplicated “0 and”.
எடுத்துக்காட்டாக, $0$ க்கு மட்டுமே அமையும் தொடர்பு
$\fn{IsZero} = \{0\}$ ஆனது பின்வரும் சொல்லால் !!{lambda definable}:
\[
  \fn{IsZero} \ident \lambd[n][n (\lambd[x][\fn{false}])\, \fn{true}].
\]
இது எவ்வாறு செயல்படுகிறது? சர்ச் எண்ணுருக்கள் மீள்செயல்படுத்திகளாக, அதாவது
தமது முதல் உள்ளீட்டை இரண்டாம் உள்ளீட்டின்மீது $n$ முறை செயல்படுத்தும்
சார்புகளாக, வரையறுக்கப்பட்டுள்ளதால் தொடக்க மதிப்பை $\fn{true}$ என
அமைக்கிறோம்; மீள்செயல்படுத்தலின் ஒவ்வொரு படியிலும், முந்தைய படியின் விளைவு
எதுவாக இருந்தாலும் $\fn{false}$ ஐத் திருப்பித்தருகிறோம். தொடக்க மதிப்பின்மீது
இப்படியை $n$ முறை செயல்படுத்துவோம். இப்படி ஒருமுறைகூடச்
செயல்படுத்தப்படாதபோது, அதாவது $n = 0$ ஆகும்போது மட்டும், விளைவு
$\fn{true}$ ஆகும்.

மெய்மதிப்புகளின் இப்பிரதிநிதித்துவத்தின் அடிப்படையில், மேலும் சில
மெய்ச்சார்புகளை வரையறுக்கலாம். மறுப்பு, இணையல் ஆகியவற்றின் இரண்டு
பிரதிநிதித்துவங்கள் இதோ:
\begin{align*}
  \fn{Not} & \ident \lambd[x][x\, \fn{false}\, \fn{true}]\\
  \fn{And} & \ident \lambd[x][\lambd[y][xy \,\fn{false}]]
\end{align*}
சார்பு “$\fn{Not}$” ஓர் உள்ளீட்டை ஏற்று, அது $\fn{false}$ எனில்
$\fn{true}$ ஐயும், அது~$\fn{true}$ எனில் $\fn{false}$ ஐயும்
திருப்பித்தருகிறது. சார்பு “$\fn{And}$” இரண்டு மெய்மதிப்புகளை உள்ளீடுகளாக
ஏற்று, இரு உள்ளீடுகளும்~$\fn{true}$ ஆக இருந்தால், அப்போதும் அப்போது மட்டுமே,
$\fn{true}$ ஐத் திருப்பித்தர வேண்டும். மெய்மதிப்புகள் மேலே விளக்கியபடி
தேர்விகளாகப் பிரதிநிதித்துவப்படுத்தப்படுகின்றன. ஆகவே $x$ ஒரு மெய்மதிப்பாக
இருக்கும்போது இரண்டு உள்ளீடுகள்மீது அதைச் செயல்படுத்தினால், $x$ ஆனது
$\fn{true}$ எனில் முதல் உள்ளீடும், இல்லையெனில் இரண்டாம் உள்ளீடும் விளைவாகும்.
இப்போது $\fn{And}$ தனது $x$, $y$ என்ற இரண்டு உள்ளீடுகளை ஏற்று, முதல்
உள்ளீடான~$x$ இடம் $y$, $\fn{false}$ ஆகியவற்றை உள்ளீடுகளாக வழங்குகிறது.
$x$ ஒரு மெய்மதிப்பு எனக் கொண்டால், $x$ ஆனது $\fn{true}$ எனில் விளைவு~$y$;
$x$ ஆனது $\fn{false}$ எனில் விளைவு $\fn{false}$. இதுவே வேண்டியது.

$\fn{And}$ க்கு மெய்மதிப்புகள் மட்டுமே உள்ளீடுகளாகப் பயன்படுத்தப்படுகின்றன
என்று இங்கு கருதுகிறோம். வேறு சொற்களை அதற்கு வழங்கினால், விளைவு (அதாவது,
இயல்புவடிவம் ஒன்று இருந்தால் அது) ஒரு மெய்மதிப்பாக இல்லாமல் போகலாம்.

\begin{prob}
  மெய்மதிப்புகளை $\lambd$-சொற்களாகக் குறியாக்கியுள்ளதைப் பயன்படுத்தி,
  உள்ளடக்கிய, விலக்கும் பிரிப்பிணைவுகளின் மெய்ச்சார்புகளைப் பிரதிநிதித்துவப்படுத்தும்
  $\fn{Or}$, $\fn{Xor}$ ஆகிய சார்புகளை வரையறுக்க.
\end{prob}

\end{document}
